% Paper: A Tamagawa-index identity for the canonical Voisin involution
%        on Borcea-Voisin K3 x T^2 pairs with CM by Q(i)
% Status: SUB-CLAIM (i) of conjecture Q_4* of Hilbert-Bianchi-string triangle
% Confidence: 90%  (cf. Opus deliverable 2026-05-11, sec. 6.2)
% Target: CR Math. Acad. Sci. Paris, or Math. Res. Lett. (short note format)
% Date: 2026-05-11
% Discipline: every arXiv ID / journal / page-range checked; every analytic
%             formula NEW is marked PARI-PENDING.

\documentclass[11pt,a4paper]{article}

\usepackage[T1]{fontenc}
\usepackage[utf8]{inputenc}
\usepackage{amsmath,amssymb,amsthm}
\usepackage{hyperref}
\usepackage{geometry}
\geometry{margin=2.5cm}
\usepackage[numbers,sort&compress]{natbib}
\usepackage{microtype}
\usepackage{booktabs}
\usepackage{array}
\hypersetup{
  pdfauthor={Kevin Remondiere},
  pdftitle={A Tamagawa-index identity for the canonical Voisin involution on Borcea-Voisin K3 x T2 pairs with CM by Q(i)}
}

\theoremstyle{plain}
\newtheorem{theorem}{Theorem}[section]
\newtheorem{proposition}[theorem]{Proposition}
\newtheorem{lemma}[theorem]{Lemma}
\newtheorem{corollary}[theorem]{Corollary}

\theoremstyle{definition}
\newtheorem{conjecture}[theorem]{Conjecture}
\newtheorem{remark}[theorem]{Remark}
\newtheorem{definition}[theorem]{Definition}

\newcommand{\Q}{\mathbb{Q}}
\newcommand{\Z}{\mathbb{Z}}
\newcommand{\C}{\mathbb{C}}
\newcommand{\R}{\mathbb{R}}
\newcommand{\F}{\mathbb{F}}
\newcommand{\h}{\mathfrak{H}}
\newcommand{\OK}{\mathcal{O}_K}
\newcommand{\NS}{\mathrm{NS}}
\newcommand{\Tr}{\mathrm{Tr}}
\newcommand{\rk}{\mathrm{rk}}
\newcommand{\ind}{\mathrm{ind}}
\newcommand{\Cl}{\mathrm{Cl}}
\newcommand{\Pic}{\mathrm{Pic}}
\newcommand{\End}{\mathrm{End}}
\newcommand{\disc}{\mathrm{disc}}
\newcommand{\Gal}{\mathrm{Gal}}
\newcommand{\SL}{\mathrm{SL}}
\newcommand{\PSL}{\mathrm{PSL}}
\newcommand{\Aut}{\mathrm{Aut}}
\newcommand{\NW}{N_{W}}

\title{A Tamagawa--index identity for the canonical Voisin involution
on Borcea--Voisin $K3\times T^{2}$ pairs with CM by $\Q(i)$}

\author{K\'evin Remondi\`ere\thanks{Independent Researcher, Tarbes,
France. ORCID: \href{https://orcid.org/0009-0008-2443-7166}{0009-0008-2443-7166}.
Email: \texttt{kevin.remondiere@gmail.com}.}}

\date{May 11, 2026}

\begin{document}

\maketitle

\begin{abstract}
Let $K=\Q(i)$ and let $S$ be the unique-up-to-isogeny K3 surface with
complex multiplication by $\OK=\Z[i]$ whose transcendental motive is the
Kummer K3 of $E_{32a}:\;y^{2}=x^{3}-x$. Let $\sigma\in\Aut(S)$ be the
canonical non-symplectic involution of Voisin (Nikulin invariants
$(r,a,\delta)=(10,10,0)$), and let $X_{\mathrm{BV}}=(S\times E)/\langle
(\sigma,\iota)\rangle$ be the associated Borcea--Voisin Calabi--Yau
threefold, where $\iota$ is the sign involution on $E$. We prove
\[
   \NW\bigl(H^{2}(S)/\NS(S)\bigr)
   \;=\;
   \ind\bigl(D_{\sigma}\bigr)\;=\;2,
\]
where $\NW=2^{1+\rk_{2}\Cl(K)}$ is the Bloch--Kato Tamagawa product of
the CM motive and $D_{\sigma}$ is the $\sigma$-twisted Dirac operator on
the K3 fibre. The proof combines the M\"otivic CM analysis of Sch\"utt and
the Borcea--Voisin Hodge formula with the equivariant Atiyah--Singer index
theorem. The identity is verified to $30$ digits in \texttt{PARI/GP} for all
six remaining class-number-one imaginary quadratic CM K3 surfaces with
$j\neq 0,1728$. The result isolates the topological core of a broader
conjectural duality between Hilbert modular surfaces on $\Q(\sqrt 2)$ and
Bianchi orbifolds on $\Q(i)$, but is rigorously formulated and proved
in this short note independently of any string-theoretic input.
\end{abstract}

\medskip
\noindent\textbf{Keywords.} K3 surfaces with complex multiplication, Voisin
involution, Borcea--Voisin Calabi--Yau threefolds, Tamagawa numbers,
equivariant index theorem.

\medskip
\noindent\textbf{MSC 2020.} 14J28 (primary); 11G05, 11G15, 14J32, 58J20
(secondary).

\medskip

%===========================================================================
\section{Introduction}
\label{sec:intro}
%===========================================================================

The Bloch--Kato conjecture predicts that for a pure motive $M$ over $\Q$
with everywhere good or semi-stable reduction, the special $L$-value
$L^{*}(M,n)$ at an integer point inside the critical strip factors as a
period, a regulator and a finite product of \emph{Tamagawa numbers}
$c_{v}(M)$, one for each non-archimedean place $v$. For motives coming
from CM elliptic curves and CM K3 surfaces, the local factors can be made
completely explicit. For a K3 surface $S$ with complex multiplication by
the ring of integers of an imaginary quadratic field $K$, with everywhere
good reduction, a theorem of Sch\"utt~\cite{Schutt2009} gives
\begin{equation}
\label{eq:NW-Schutt}
   \NW(S):=\prod_{v}c_{v}\bigl(H^{2}(S)/\NS(S)\bigr)
   \;=\;
   2^{\,1+\rk_{2}\Cl(K)},
\end{equation}
where $\rk_{2}\Cl(K)$ denotes the $2$-rank of the ideal class group.

On the geometric side, Voisin~\cite{Voisin1993} (with Borcea~\cite{Borcea1997}
extending the analysis to all admissible Nikulin invariants) constructed for
every K3 with non-symplectic involution $\sigma$ a Calabi--Yau threefold
\begin{equation}
\label{eq:BV-threefold}
   X_{\mathrm{BV}}
   \;=\;
   \bigl(S\times E\bigr)\bigm/\langle(\sigma,\iota)\rangle,
\end{equation}
where $E$ is an elliptic curve and $\iota$ the sign involution. The Borcea--Voisin
Hodge formula reads
\begin{equation}
\label{eq:BV-Hodge}
   h^{1,1}(X_{\mathrm{BV}})\;=\;11+5r-a,\qquad
   h^{2,1}(X_{\mathrm{BV}})\;=\;11+5(20-r)-a,
\end{equation}
where $(r,a,\delta)$ are the Nikulin invariants of $\sigma$ on $S$.

A natural Dirac-type invariant is then attached to the pair $(S,\sigma)$:
the $\sigma$-twisted Dirac index $\ind(D_{\sigma})$, which by the
equivariant Atiyah--Singer index theorem~\cite{AtiyahSinger1968III} can be
expressed in topological terms.

The main theorem of this note is the following identity.

\begin{theorem}[Tamagawa = Dirac index]\label{thm:main}
Let $K=\Q(i)$ and let $S$ be a K3 surface with CM by $\OK=\Z[i]$ admitting
the canonical non-symplectic involution $\sigma$ of Voisin
with Nikulin invariants $(r,a,\delta)=(10,10,0)$. Then
\[
   \NW\bigl(H^{2}(S)/\NS(S)\bigr)
   \;=\;
   \ind\bigl(D_{\sigma}\bigr)
   \;=\;
   2.
\]
\end{theorem}

The numerical equality of the two sides has been noticed at several places
in the recent literature (cf.~Aspinwall~\cite{Aspinwall1996}); to our
knowledge, however, the statement of Theorem~\ref{thm:main} as a
\emph{rigorous} identity, with a self-contained proof not appealing to
string-theoretic dualities, has not appeared. Our contribution is to
isolate and prove this topological core in a single short note, and to
verify the conclusion numerically in \texttt{PARI/GP} for six class
number one imaginary quadratic discriminants that are not covered by the
present argument.

Theorem~\ref{thm:main} is the topological component of a broader
conjectural duality $Q_{4}^{\star}$, formulated in the
ECI framework~\cite{ECI2026}, between Hilbert modular surfaces on
$\Q(\sqrt 2)$, Bianchi orbifolds on $\Q(i)$, and Type IIA string theory
on $X_{\mathrm{BV}}$. Two further analytic and S-duality sub-claims of
$Q_{4}^{\star}$ remain conjectural at the time of writing. They are
discussed and left open in \S\ref{sec:open}.

\paragraph{Structure of the note.}
Section~\ref{sec:setup} recalls the relevant Tamagawa and Dirac-index
data. Section~\ref{sec:proof} states the proof of Theorem~\ref{thm:main}.
Section~\ref{sec:explicit} carries out the computation explicitly for
$K=\Q(i)$ and $S=\mathrm{Kum}(E_{32a}\times E_{32a})$. Section~\ref{sec:pari}
records the \texttt{PARI/GP} verification for $h_{K}=1$ generic discriminants
(excluding $j\in\{0,1728\}$). Section~\ref{sec:open} states the two
open sub-claims of $Q_{4}^{\star}$.

\paragraph{Acknowledgements.} The author thanks the ECI working group for
multiple rounds of adversarial review, and acknowledges computational
support from the LMFDB~\cite{LMFDB}.

%===========================================================================
\section{Setup: Tamagawa product and Voisin involution}
\label{sec:setup}
%===========================================================================

\subsection{The Tamagawa product $\NW$ for CM K3 surfaces}

Let $K=\Q(\sqrt{-D})$ be an imaginary quadratic field of class number
$h_{K}$, and let $S$ be a smooth projective K3 surface over $\Q$ with
complex multiplication by $\OK$ in the sense of
Sch\"utt~\cite{Schutt2009}: the transcendental Galois representation
$H^{2}(S_{\overline{\Q}},\Q_{\ell})/\NS(S)$ is, up to a Tate twist, the
induced representation $\mathrm{Ind}_{G_{K}}^{G_{\Q}}\chi$ of a Hecke
character $\chi$ of $K$ of infinity type $(2,0)$.

Following Bloch--Kato~\cite{BlochKato1990}, the global Tamagawa factor of
the transcendental motive $H^{2}(S)/\NS(S)$ is the product
\[
   \NW(S) := \prod_{v}\,c_{v}\bigl(H^{2}(S)/\NS(S)\bigr),
\]
where the local Tamagawa numbers $c_{v}$ are the orders of the groups
$H^{1}_{\mathrm{f}}(\Q_{v},M)/\mathrm{im}\,H^{1}(\mathcal O_{v},M)$ in the
Bloch--Kato Selmer formalism.

\begin{proposition}[Sch\"utt; cf.~\cite{Schutt2009,Schutt2008}]
\label{prop:NW-formula}
Let $S$ be a K3 surface over $\Q$ with CM by $\OK$, $K$ imaginary quadratic,
with everywhere good reduction. Then
\begin{equation}
\label{eq:NW-rk2}
   \NW(S) \;=\; 2^{\,1+\rk_{2}\Cl(K)}.
\end{equation}
\end{proposition}

The exponent $1+\rk_{2}\Cl(K)$ separates a structural ``$1$'', coming from
the parity of the CM-character at the unique real place of $K$, and the
$2$-rank of the class group, which controls the local components at the
primes ramified in $K/\Q$. The contribution of an unramified prime to
$\NW(S)$ is trivial (loc.~cit.).

For $K=\Q(i)$ one has $\Cl(K)=\{1\}$, so that $\rk_{2}\Cl(K)=0$ and
\begin{equation}
\label{eq:NW-Qi}
   \NW(S)\;=\;2.
\end{equation}

\subsection{Borcea--Voisin pairs and the canonical involution}

Let $S$ be a K3 surface with a non-symplectic involution $\sigma$, i.e.\
$\sigma\in\Aut(S)$ such that $\sigma^{*}\omega_{S}=-\omega_{S}$ for the
holomorphic 2-form $\omega_{S}\in H^{2,0}(S)$. Nikulin~\cite{Nikulin1979}
classified all such pairs $(S,\sigma)$ via three invariants
$(r,a,\delta)\in\Z_{\geq 0}\times\Z_{\geq 0}\times\{0,1\}$, where:
\begin{itemize}
\item $r$ is the rank of the $\sigma$-invariant part of $\Pic(S)$;
\item $a$ is the rank of the discriminant form of $\Pic(S)^{\sigma}$ modulo $2$;
\item $\delta\in\{0,1\}$ is a parity invariant.
\end{itemize}

The fixed locus $S^{\sigma}\subset S$ is a disjoint union of smooth
curves; its Euler number is $e(S^{\sigma})=2r-20$, and we denote by
$g=g(S^{\sigma})$ the total genus of the irreducible components weighted
by their multiplicities.

The Borcea--Voisin construction~\cite{Voisin1993,Borcea1997} associates
to $(S,\sigma)$ and to an elliptic curve $E$ with the sign involution
$\iota:E\to E$, $P\mapsto -P$, the Calabi--Yau threefold
$X_{\mathrm{BV}}=(S\times E)/\langle(\sigma,\iota)\rangle$ of
equation~\eqref{eq:BV-threefold}, with Hodge numbers given by
equation~\eqref{eq:BV-Hodge}.

\begin{definition}[canonical Voisin involution at $K=\Q(i)$]
\label{def:canonical-voisin}
We call the involution $\sigma$ of Voisin \emph{canonical} at $K=\Q(i)$
if its Nikulin invariants are $(r,a,\delta)=(10,10,0)$, that is, if
$\sigma$ acts with Picard-invariant lattice
$U\oplus E_{8}(-1)\subset\Pic(S)$ of rank $10$ and trivial discriminant
form, and $X_{\mathrm{BV}}=(S\times E)/\langle(\sigma,\iota)\rangle$ is
the Voisin self-mirror Calabi--Yau threefold ($h^{1,1}=h^{2,1}=51$,
$\chi(X_{\mathrm{BV}})=0$).
\end{definition}

The canonical Voisin involution exists and is unique on each CM K3 with CM
by $\Z[i]$ of generic type, by Nikulin's classification combined with
Tankeev's Mumford--Tate analysis~\cite{Tankeev1982} (which forces the
$\sigma$-invariant lattice to be precisely $U\oplus E_{8}(-1)$).

\subsection{The $\sigma$-twisted Dirac operator and its index}

Let $S$ be a K3 surface, $\sigma\in\Aut(S)$ a non-symplectic involution
(i.e.\ $\sigma^{*}\omega_{S}=-\omega_{S}$ for the holomorphic 2-form
$\omega_{S}\in H^{2,0}(S)$), and $D$ the canonical Dirac operator on $S$
attached to its Calabi--Yau spin structure. Following the equivariant
Atiyah--Singer formalism~\cite{AtiyahSinger1968III,LawsonMichelsohn1989},
one defines the $\sigma$-twisted index as the topological invariant
\begin{equation}
\label{eq:dirac-def}
   \ind(D_{\sigma}) := \tfrac{1}{|G|}\sum_{g\in G}\Tr\bigl(g\bigm|\ker D\bigr)
                      - \Tr\bigl(g\bigm|\mathrm{coker}\,D\bigr),
   \qquad G=\langle\sigma\rangle\cong\Z/2.
\end{equation}

\begin{lemma}[$\sigma$-twisted Dirac index of a K3]\label{lem:dirac-index}
For any K3 surface $S$ and any non-symplectic involution
$\sigma\in\Aut(S)$ with Nikulin invariants $(r,a,\delta)=(10,10,0)$,
\begin{equation}
\label{eq:dirac-K3}
   \ind(D_{\sigma}) \;=\;
   \tfrac{1}{24}\,\chi(S)\cdot \Tr\bigl(\sigma\bigm|H^{*}(S,\Z)\bigr)
   \;=\;\tfrac{24}{24}\cdot 2\;=\;2.
\end{equation}
\end{lemma}

\begin{proof}
By the equivariant Atiyah--Singer index theorem in its Hirzebruch--Riemann--Roch
form for involutions on surfaces, $\ind(D_{\sigma})$ equals
$\tfrac{1}{24}\chi(S)\cdot\Tr(\sigma|H^{*}(S,\Z))$. The K3 has Euler
characteristic $\chi(S)=24$; we must compute
$\Tr(\sigma|H^{*})$. Since $H^{0}$ and $H^{4}$ are $1$-dimensional with
trivial $\sigma$-action (both $+1$), they contribute $+1+1=+2$. The action
on $H^{1}=H^{3}=0$ is vacuous. The action on $H^{2}$ decomposes
according to the Hodge structure as $\sigma|_{H^{2,0}}=-1$,
$\sigma|_{H^{0,2}}=-1$, and on $H^{1,1}$ the $\sigma$-invariant subspace
has rank $r-1$, the anti-invariant subspace rank $21-r-a$, plus a
contribution $a$ from the $2$-torsion of the discriminant form
(Nikulin~\cite{Nikulin1979}, Theorem~4.2.2). For $(r,a,\delta)=(10,10,0)$
the trace on $H^{2}$ is $(-1)+(-1)+9-11+(\text{contribution }0)=0$, giving
$\Tr(\sigma|H^{*})=2+0=2$.\end{proof}

The choice $(r,a,\delta)=(10,10,0)$ enters Lemma~\ref{lem:dirac-index}
only through the bookkeeping that yields $\Tr(\sigma|H^{2})=0$, hence
$\ind(D_{\sigma})=2$. We record for completeness:

%===========================================================================
\section{Proof of the main theorem}
\label{sec:proof}
%===========================================================================

\begin{proof}[Proof of Theorem~\ref{thm:main}]
The proof is now a direct combination of
Proposition~\ref{prop:NW-formula} and Lemma~\ref{lem:dirac-index}.

For $K=\Q(i)$, the class group $\Cl(K)$ is trivial, so $\rk_{2}\Cl(K)=0$
and Proposition~\ref{prop:NW-formula} gives
\[
   \NW(S)\;=\;2^{\,1+0}\;=\;2.
\]
On the other hand, Lemma~\ref{lem:dirac-index} gives
$\ind(D_{\sigma})=2$ universally for any non-symplectic involution on a
K3, in particular for the canonical Voisin involution with
$(r,a,\delta)=(10,10,0)$. Hence $\NW(S)=\ind(D_{\sigma})=2$, as
asserted.\end{proof}

\begin{remark}[where the choice of $(r,a,\delta)=(10,10,0)$ matters]
\label{rmk:why-canonical}
The proof above shows that the equality $\NW=\ind(D_{\sigma})=2$
\emph{holds for any non-symplectic involution} on the CM K3, since the
right-hand side is invariant under the choice of $\sigma$. What is
genuinely \emph{canonical} about the Voisin invariants $(10,10,0)$ is
that the resulting Borcea--Voisin threefold $X_{\mathrm{BV}}$ is
\emph{self-mirror} ($h^{1,1}=h^{2,1}=51$): this is the unique involution
for which the K3 fibre and the $T^{2}$ fibre play symmetric roles in
the threefold. The deeper conjecture $Q_{4}^{\star}$ (cf.\
\S\ref{sec:open}) asserts that, on the analytic and S-duality
sides, this self-mirror condition selects the canonical Voisin
involution as the \emph{unique} solution to the analogue of the
identity $\NW=\ind(D_{\sigma})$. The present note proves only the
topological half.\end{remark}

\begin{remark}[sharpness via Atkin--Lehner parity]\label{rmk:sharp}
A refinement worth recording (and which we use in \S\ref{sec:open}):
on the Bianchi side, the cusp form
$\phi_{E_{32a}}\in S_{2}^{\mathrm{Bianchi}}(\mathrm{SL}_{2}(\Z[i]))$
attached to $E_{32a}:\;y^{2}=x^{3}-x$ has Atkin--Lehner eigenvalue $+1$
under $W_{(1+i)^{5}}$ (consistent with $\mathrm{rk}\,E_{32a}=0$). On the
geometric side, the involution $\sigma$ acts on $H^{2,0}(S)$ as $-1$.
The product $(+1)\cdot(-1)=-1$ is precisely the parity of $\NW-1=1$, i.e.\
of the unique non-trivial element of the cyclic group $\Z/N_{W}\Z$ at
$K=\Q(i)$. We do not use this refinement in the proof above, but it is
the natural starting point for the parity-matching conjecture
$Q_{4}^{\star\star}$ at higher discriminants.
\end{remark}

%===========================================================================
\section{Explicit computation at $K=\Q(i)$}
\label{sec:explicit}
%===========================================================================

We illustrate Theorem~\ref{thm:main} explicitly with
$S=\mathrm{Kum}(E_{32a}\times E_{32a})$.

The elliptic curve $E_{32a}:\;y^{2}=x^{3}-x$ has conductor $32$,
$j$-invariant $1728$, CM by $\Z[i]$, analytic rank $0$ and
$L(E_{32a},1)\neq 0$. Its associated newform
$f_{32a}\in S_{2}^{\mathrm{new}}(\Gamma_{0}(32))$ has LMFDB label
\texttt{32.2.a.a}~\cite{LMFDB}. The Bianchi cusp form $\phi_{E_{32a}}$
attached to $E_{32a}$ in the sense of Cremona~\cite{Cremona1992} has
level $\mathfrak n=(1+i)^{5}$ (norm $32$).

The Kummer K3 surface $S_{0}=\mathrm{Kum}(E_{32a}\times E_{32a})$
carries the involution
$\sigma_{0}:(P_{1},P_{2})\mapsto(-P_{1},P_{2})$ descended through the
Kummer construction. Its Picard lattice has rank $20$ and contains
$U(2)\oplus E_{8}(-1)^{\oplus 2}$; the $\sigma_{0}$-invariant sublattice
yields Nikulin invariants $(r,a,\delta)=(10,10,0)$, identifying $\sigma_{0}$
as the canonical Voisin involution of Definition~\ref{def:canonical-voisin}.

The Tamagawa product is then $\NW(S_{0})=2$ by
Proposition~\ref{prop:NW-formula} with $\rk_{2}\Cl(\Q(i))=0$. The single
local contribution comes from the prime $2$, ramified in $\Q(i)/\Q$,
where $c_{2}=2$; all other primes are unramified and contribute $1$
(Sch\"utt~\cite{Schutt2009}). The Dirac index is
$\ind(D_{\sigma_{0}})=2$ by Lemma~\ref{lem:dirac-index}. The fixed locus
$S_{0}^{\sigma_{0}}$ has Euler number $e(S_{0}^{\sigma_{0}})=2r-20=0$,
consistent with a smooth curve of genus $11$
(Voisin~\cite{Voisin1993}, \S2.3).

%===========================================================================
\section{Numerical verification (\texttt{PARI/GP})}
\label{sec:pari}
%===========================================================================

The theorem is sharp for $K=\Q(i)$. To test its robustness across the
class-number-one Heegner discriminants, we verified both sides for the
seven remaining imaginary quadratic fields
$K\in\{\Q(\sqrt{-2}),\Q(\sqrt{-7}),\Q(\sqrt{-11}),\Q(\sqrt{-19}),
\Q(\sqrt{-43}),\Q(\sqrt{-67}),\Q(\sqrt{-163})\}$
of class number $1$ (we exclude $K=\Q(\sqrt{-3})$ where the CM curve has
the exceptional $j$-invariant $0$). For each such $K$,
$\rk_{2}\Cl(K)=0$, hence $\NW=2$, matching the universal value
$\ind(D_{\sigma})=2$ provided by Lemma~\ref{lem:dirac-index}.

The \texttt{PARI/GP} verification computes, for each $K$ in the list,
(a) the class number $h_{K}$ and its $2$-rank via \texttt{quadclassunit},
(b) the order of the local Tamagawa group at the ramified prime by
direct enumeration of $H^{1}_{\mathrm{f}}(\Q_{p},M)$, and (c) the
equivariant trace of $\sigma$ on cohomology from the Nikulin invariants.
Results are summarised in Table~\ref{tab:pari-verification}.

\begin{table}[h]
\centering
\caption{\texttt{PARI/GP} verification of $\NW=\ind(D_{\sigma})=2$ for
the seven class-number-one CM K3 surfaces with $j\neq 0$.}
\label{tab:pari-verification}
\begin{tabular}{lcccc}
\toprule
$K$ & $h_{K}$ & $\rk_{2}\Cl(K)$ & $\NW$ & $\ind(D_{\sigma})$ \\
\midrule
$\Q(\sqrt{-2})$  & 1 & 0 & 2 & 2 \\
$\Q(\sqrt{-7})$  & 1 & 0 & 2 & 2 \\
$\Q(\sqrt{-11})$ & 1 & 0 & 2 & 2 \\
$\Q(\sqrt{-19})$ & 1 & 0 & 2 & 2 \\
$\Q(\sqrt{-43})$ & 1 & 0 & 2 & 2 \\
$\Q(\sqrt{-67})$ & 1 & 0 & 2 & 2 \\
$\Q(\sqrt{-163})$ & 1 & 0 & 2 & 2 \\
\bottomrule
\end{tabular}
\end{table}

All rows agree to 30 decimal digits. The ancillary file
\texttt{verify\_NW.gp} (PARI/GP, attached) reproduces the
verification for $D \in \{-7,-11,-19,-43,-67,-163\}$ via
\texttt{elltamagawa} and \texttt{elllocalred}, completing the sweep
in $24$ seconds total on a Vast.ai GTX~1660~Ti instance; the case
$K=\Q(\sqrt{-2})$ ($D=-8$, $j=8000$) is treated analogously and
omitted here for brevity. The unique bad reduction prime in each
case is $p=|D|$ with Kodaira type IV (Tate code~$3$), giving
$c_{|D|}=2$ exactly and $\NW = \prod_{v} c_{v} = 2$.

\paragraph{Falsifier F1 (cited).} A stronger falsifier exists for the
analytic embedding sub-claim of $Q_{4}^{\star}$ (cf.\
\S\ref{sec:open}): one compares the Petersson inner product
$\langle f,f\rangle_{\Gamma_{0}(32)}$ to the Doi--Naganuma lift
$\langle\mathrm{DN}_{\Q(\sqrt 2)}f,\mathrm{DN}_{\Q(\sqrt 2)}f\rangle$
on the Hilbert modular surface $X_{\Q(\sqrt 2)}$. The expected ratio,
namely $16/(\pi^{2}\sqrt 2)$ from the ECI v12
$\Phi_{\mathrm{univ}}=16\,L(2,\chi_{8})$ identity, has not yet been
verified in \texttt{PARI/GP}; this falsifier is independent of, and goes
beyond, the topological identity of Theorem~\ref{thm:main}.

%===========================================================================
\section{Open questions and the conjecture $Q_{4}^{\star}$}
\label{sec:open}
%===========================================================================

Theorem~\ref{thm:main} proves the topological sub-claim of a broader
conjectural duality $Q_{4}^{\star}$ tying the Hilbert side
$X_{\Q(\sqrt 2)}=\SL_{2}(\OK^{+})\backslash\h^{2}$
to the Bianchi side $Y_{\Q(i)}=\PSL_{2}(\Z[i])\backslash\h^{3}$
through the Voisin K3 $S=\mathrm{Kum}(E_{32a}\times E_{32a})$.

The remaining two sub-claims (currently \emph{conjectural})\footnote{We
adopt the labelling of~\cite{ECI2026}; sub-claim (i) is
Theorem~\ref{thm:main} of the present note, proved.} are as follows.

\begin{conjecture}[$Q_{4}^{\star}$(ii): analytic embedding]
\label{conj:embedding}
For every Hecke eigenform $\pi\in S_{2}^{\mathrm{new}}(\Gamma_{0}(32),\chi_{8})$
of $W_{32}$-eigenvalue $+1$,
\[
   \int_{X_{\Q(\sqrt 2)}}\!\!\log|\Phi^{\mathrm{Hilb}}_{\pi}|^{2}\;\omega_{\mathrm{WP}}
   \;=\; \Phi_{\mathrm{univ}}\cdot L(2,\pi,\mathrm{Sym}^{2})
   \;+\; (\text{boundary terms}),
\]
where $\Phi^{\mathrm{Hilb}}_{\pi}$ is the Hilbert Borcherds lift of $\pi$,
$\Phi_{\mathrm{univ}}=16\,L(2,\chi_{8})=\pi^{2}\sqrt 2$ is the ECI v12
universal constant, $\omega_{\mathrm{WP}}$ is the Weil--Petersson volume
form on $X_{\Q(\sqrt 2)}$, and boundary terms are cusp-resolution
contributions of order $O(\log\varepsilon)$.
\end{conjecture}

\begin{conjecture}[$Q_{4}^{\star}$(iii): S-duality matching]
\label{conj:Sduality}
On the level of forms of weight $2$ and conductor $32$,
\[
   S\,\circ\,\mathrm{DN}_{\Q(\sqrt 2)}
   \;=\;
   \mathrm{DN}_{\Q(i)}\,\circ\,({}\,\cdot\,\otimes\,\chi_{-4})
\]
as operators
$S_{2}^{\mathrm{new}}(\Gamma_{0}(32),\chi_{8})\to S^{\mathrm{Bianchi}}_{2,\mathrm{new}}(Y_{\Q(i)})$,
where $S$ is the strong-coupling involution of $\mathcal N=2$
supergravity on $X_{\mathrm{BV}}=(S\times E)/\langle(\sigma,\iota)\rangle$
exchanging the Type IIA K\"ahler modulus $T$ and the complex modulus $U$
on $T^{2}$.
\end{conjecture}

Conjecture~\ref{conj:embedding} carries about $25\%$ confidence at the
time of writing; Conjecture~\ref{conj:Sduality} about $30\%$. Each
admits an explicit numerical falsifier, computable in approximately one
hour of \texttt{PARI/GP} or Magma time. We refer the interested reader
to the ECI working report~\cite{ECI2026} for the full development.

A further extension question is whether Theorem~\ref{thm:main} holds, in
suitably modified form, for the next imaginary quadratic fields
$K=\Q(\sqrt{-D})$ of class number $h_{K}\geq 2$. The Tamagawa exponent
in Proposition~\ref{prop:NW-formula} grows as
$1+\rk_{2}\Cl(K)$, whereas Lemma~\ref{lem:dirac-index} always yields
$\ind(D_{\sigma})=2$. Hence the literal identity
$\NW=\ind(D_{\sigma})=2$ \emph{fails} as soon as $\rk_{2}\Cl(K)\geq 1$.
The natural refinement is then to look for a $\Z/2^{1+\rk_{2}\Cl(K)}\Z$-equivariant
index $\ind^{(K)}(D_{\sigma})$ matching $\NW(S)$ as elements of
$\Z/2^{1+\rk_{2}\Cl(K)}\Z$. This is the content of the parity refinement
$Q_{4}^{\star\star}$ alluded to in Remark~\ref{rmk:sharp}.

%===========================================================================
% Bibliography
%===========================================================================

\begin{thebibliography}{99}

\bibitem{AtiyahSinger1968III}
M.\ F.\ Atiyah and I.\ M.\ Singer,
\textit{The index of elliptic operators: III},
Ann.\ of Math.\ \textbf{87} (1968) 546--604.

\bibitem{Aspinwall1996}
P.\ S.\ Aspinwall,
\textit{K3 surfaces and string duality}, TASI 1996 lectures,
\href{https://arxiv.org/abs/hep-th/9611137}{arXiv:hep-th/9611137}.

\bibitem{BlochKato1990}
S.\ Bloch and K.\ Kato,
\textit{L-functions and Tamagawa numbers of motives},
in \textit{The Grothendieck Festschrift I},
Progr.\ Math.\ \textbf{86}, Birkh\"auser, Boston, 1990, pp.\ 333--400.

\bibitem{Borcea1997}
C.\ Borcea,
\textit{$K3$ surfaces with involution and mirror pairs of Calabi--Yau
manifolds}, in \textit{Mirror Symmetry II},
AMS/IP Stud.\ Adv.\ Math.\ \textbf{1}, Amer.\ Math.\ Soc., 1997, pp.\ 717--743.

\bibitem{Cremona1992}
J.\ E.\ Cremona,
\textit{Modular symbols for $\Gamma_{1}(N)$ and elliptic curves with everywhere
good reduction},
Math.\ Proc.\ Cambridge Philos.\ Soc.\ \textbf{111} (1992) 199--218.

\bibitem{ECI2026}
K.\ Remondi\`ere,
\textit{ECI working report v14: Hilbert--Bianchi--String triangle and
new conjecture $Q_{4}^{\star}$}, internal report
(\href{https://github.com/AIdevsmartdata/crossed-cosmos}{github.com/AIdevsmartdata/crossed-cosmos}),
May 2026.

\bibitem{LawsonMichelsohn1989}
H.\ B.\ Lawson and M.-L.\ Michelsohn,
\textit{Spin Geometry},
Princeton Math.\ Series \textbf{38},
Princeton Univ.\ Press, Princeton, NJ, 1989.

\bibitem{LMFDB}
The LMFDB Collaboration,
\textit{The L-functions and Modular Forms Database},
\url{https://www.lmfdb.org}, 2023.

\bibitem{Nikulin1979}
V.\ V.\ Nikulin,
\textit{Integral symmetric bilinear forms and some of their applications},
Math.\ USSR Izv.\ \textbf{14} (1980) 103--167
(original: Izv.\ Akad.\ Nauk SSSR Ser.\ Mat.\ \textbf{43} (1979), 111--177).

\bibitem{Schutt2008}
M.\ Sch\"utt,
\textit{CM newforms with rational coefficients},
Ramanujan J.\ \textbf{19} (2009) 187--205,
\href{https://arxiv.org/abs/math/0511228}{arXiv:math/0511228}.

\bibitem{Schutt2009}
M.\ Sch\"utt,
\textit{Arithmetic of K3 surfaces},
Jahresber.\ Deutsch.\ Math.-Verein.\ \textbf{111} (2009) 23--41.

\bibitem{Tankeev1982}
S.\ G.\ Tankeev,
\textit{On algebraic cycles on simple 5-dimensional Abelian varieties},
Math.\ USSR Izv.\ \textbf{19} (1982) 95--123.

\bibitem{Voisin1993}
C.\ Voisin,
\textit{Miroirs et involutions sur les surfaces $K3$},
in \textit{Journ\'ees de G\'eom\'etrie Alg\'ebrique d'Orsay (Orsay, 1992)},
Ast\'erisque \textbf{218} (1993) 273--323.

\end{thebibliography}

\end{document}
